\section{Kết luận và Hướng Phát triển}

Bài báo này đã trình bày một kiến trúc ghép cặp hai chiều hoàn toàn không cần huấn luyện (training-free) và tối ưu cho phần cứng CPU cho bài toán hồ sơ ứng viên và mô tả công việc, dựa trên biểu diễn đa vector nhận biết phân mục (Section-Aware Multi-Vector) kết hợp động cơ xếp hạng lai (Calibrated Hybrid Engine) và cây tri thức kỹ năng 3 tầng. Bằng cách bóc tách văn bản thành các đơn vị chức năng độc lập và áp dụng cơ chế ghép cặp có tái chuẩn hóa trọng số động, phương pháp đã loại bỏ triệt để hiện tượng cắt cụt văn bản ($>95\%$ CV tránh truncation), khắc phục lỗi dính chữ tiêu đề thực tế, và cung cấp bảng truy vết điểm số minh bạch đến từng phân mục và quan hệ phả hệ kỹ năng.

Thực nghiệm diện rộng trên hai tập chuẩn quốc tế (\textsc{resume-job-description-fit} với $1.759$ cặp kiểm thử, $76$ truy vấn huấn luyện sạch tách biệt, và tập mật độ cao \textsc{med2425} với $8.575$ cặp) cùng các kiểm định thống kê chặt chẽ đã mang lại ba kết luận khoa học cốt lõi:
\begin{enumerate}
    \item \textbf{Cơ chế Khắc phục Cắt cụt Văn bản theo Độ dài Hồ sơ}: Mặc dù đạt sự tương đương thống kê trên điểm trung bình gộp của tập nhỏ, phân tích phân tầng theo độ dài CV chỉ ra quỹ đạo tăng trưởng đơn điệu rõ rệt ($\Delta\text{nDCG@10}$ chuyển từ $-0.0112$ ở CV ngắn $\le 350$ từ lên $+0.0126$ ở CV $>350$ từ và đạt $\mathbf{+0.0239}$ ở CV rất dài $>600$ từ). Hơn nữa, ở chiều gợi ý việc làm ngược ($\text{CV} \rightarrow \text{JD}$, $N=272$ truy vấn ứng viên), biểu diễn phân mục kết hợp vượt qua mô hình phẳng có ý nghĩa thống kê ($\bar{\Delta} = +0.0185$ nDCG@10, $t=3.09, p=0.0022$).
    \item \textbf{Sức mạnh Quyết định của Cây Phân loại trong Môi trường Mật độ Cao Hai chiều}: Trong các nhóm ứng viên nhỏ ($N=25$), cây kỹ năng chủ yếu đóng vai trò kiểm toán điều kiện tiên quyết và giải trình. Tuy nhiên, khi mật độ ứng viên tăng lên quy mô thực tế ($\approx 209$ ứng viên/JD trên \textsc{med2425}), Calibrated Hybrid tạo bước nhảy vọt có ý nghĩa thống kê so với Single-Vector ở cả hai chiều sàng lọc ứng viên ($\text{JD} \rightarrow \text{CV}$: $\bar{\Delta} = \mathbf{+0.0849}$ nDCG@10, $+15.7\%$, $p=0.0211$) và gợi ý việc làm ($\text{CV} \rightarrow \text{JD}$: $\bar{\Delta} = \mathbf{+0.0855}$ nDCG@10, $+15.4\%$, $p<0.0001$), đồng thời vượt có ý nghĩa so với Dense+BM25 ($p=0.0457$).
    \item \textbf{Biên Pareto Tối ưu trên CPU so với Cross-Encoder Nơ-ron}: Với độ trễ xếp hạng chỉ $T_{\text{rank}} = 2.69$\,ms ($T_{\text{online}} = 33.44$\,ms toàn trình bao gồm mã hóa truy vấn) trên CPU thông thường, hệ thống vượt trội có ý nghĩa thống kê so với các mô hình Cross-Encoder zero-shot cồng kềnh ($\bar{\Delta} = +0.0632$ nDCG@10, $p_{\text{RQ}} = 0.0082$) với tốc độ nhanh hơn gấp $35$ lần ($1.178,05$\,ms).
\end{enumerate}

Trong tương lai, chúng tôi dự kiến tích hợp mô hình phân tích bố cục thị giác (Vision Document Layout Analysis) cho các định dạng PDF đồ họa đa cột, mở rộng cây kỹ năng theo khung chuẩn quốc tế (ESCO/O*NET) sang các ngành nghề phi kỹ thuật, và thực hiện thử nghiệm người dùng (Recruiter User Study) để định lượng thời gian thẩm định hồ sơ tiết kiệm được trong thực tế.
